\section*{3.3 长度与距离}
\addcontentsline{toc}{section}{3.3 长度与距离}
\label{sec:3.3}

在第 3.1 节中，我们已经讨论了可用于计算向量长度的范数。内积与范数之间密切相关，因为任何内积都会以自然的方式诱导出范数：

\begin{equation}
\|\boldsymbol{x}\| := \sqrt{\langle \boldsymbol{x}, \boldsymbol{x} \rangle}
\tag{3.16}\label{eq:3-16}
\end{equation}

从而使我们能够使用内积来计算向量的长度。然而，并非每个范数都是由内积诱导的。曼哈顿范数 (3.3) 就是一个没有对应内积的范数示例。在下文中，我们将重点关注由内积诱导的范数，并引入几何概念，如长度、距离和夹角。

\emph{备注} （柯西-施瓦茨不等式）。对于内积向量空间 $ (V, \langle \cdot, \cdot \rangle) $，诱导范数 $ \|\cdot\| $ 满足 \emph{柯西-施瓦茨不等式（Cauchy-Schwarz inequality）} ：

\begin{equation}
|\langle \boldsymbol{x}, \boldsymbol{y} \rangle| \leqslant \|\boldsymbol{x}\| \|\boldsymbol{y}\|
\tag{3.17}\label{eq:3-17}
\end{equation}

♦

\begin{notebox}
\textbf{例 3.5（使用内积计算向量的长度）}
\end{notebox}

\begin{notebox}
在几何学中，我们通常对向量的长度感兴趣。现在我们可以使用内积通过式~\eqref{eq:3-16} 来计算它们。取 $ \boldsymbol{x} = [1, 1]^\top \in \mathbb{R}^2 $。如果我们使用点积作为内积，根据式~\eqref{eq:3-16}，我们得到 $ \boldsymbol{x} $ 的长度为：

\begin{equation}
\|\boldsymbol{x}\| = \sqrt{\boldsymbol{x}^\top \boldsymbol{x}} = \sqrt{1^2 + 1^2} = \sqrt{2}
\tag{3.18}\label{eq:3-18}
\end{equation}

现在让我们选择另一个不同的内积：

\begin{equation}
\langle \boldsymbol{x}, \boldsymbol{y} \rangle := \boldsymbol{x}^\top \begin{bmatrix} 1 & -\frac{1}{2} \\ -\frac{1}{2} & 1 \end{bmatrix} \boldsymbol{y} = x_1 y_1 - \frac{1}{2}(x_1 y_2 + x_2 y_1) + x_2 y_2
\tag{3.19}\label{eq:3-19}
\end{equation}

如果我们计算一个向量的范数，当 $ x_1 $ 和 $ x_2 $ 同号（且 $ x_1 x_2 > 0 $）时，该内积返回的值小于点积；否则，它返回的值大于点积。在这个内积下，我们得到：

\begin{equation}
\langle \boldsymbol{x}, \boldsymbol{x} \rangle = x_1^2 - x_1 x_2 + x_2^2 = 1 - 1 + 1 = 1 \implies \|\boldsymbol{x}\| = \sqrt{1} = 1
\tag{3.20}\label{eq:3-20}
\end{equation}

因此在这个内积下，$ \boldsymbol{x} $ 比在点积下更“短”。
\end{notebox}

\textbf{定义 3.6}（距离与度量）。考虑一个内积空间 $ (V, \langle \cdot, \cdot \rangle) $。那么对于 $ \boldsymbol{x}, \boldsymbol{y} \in V $，

\begin{equation}
d(\boldsymbol{x}, \boldsymbol{y}) := \|\boldsymbol{x} - \boldsymbol{y}\| = \sqrt{\langle \boldsymbol{x} - \boldsymbol{y}, \boldsymbol{x} - \boldsymbol{y} \rangle}
\tag{3.21}\label{eq:3-21}
\end{equation}

称为 $ \boldsymbol{x} $ 与 $ \boldsymbol{y} $ 之间的 \emph{距离（distance）} 。如果我们使用点积作为内积，则该距离称为 \emph{欧几里得距离（Euclidean distance）} 。 映射：

\begin{equation}
d : V \times V \to \mathbb{R}
\tag{3.22}\label{eq:3-22}
\end{equation}

\begin{equation}
(\boldsymbol{x}, \boldsymbol{y}) \mapsto d(\boldsymbol{x}, \boldsymbol{y})
\tag{3.23}\label{eq:3-23}
\end{equation}

称为 \emph{度量（metric）} 。

\emph{备注} 。与向量的长度类似，向量之间的距离并不一定需要内积：范数就足够了。如果我们使用的是由内积诱导的范数，则距离可能会随着内积的选择而变化。♦

度量 $ d $ 满足以下性质：

\begin{enumerate}
\item $ d $ 是 \emph{正定的（positive definite）} ，即对所有 $ \boldsymbol{x}, \boldsymbol{y} \in V $，有 $ d(\boldsymbol{x}, \boldsymbol{y}) \geqslant 0 $，且 $ d(\boldsymbol{x}, \boldsymbol{y}) = 0 \iff \boldsymbol{x} = \boldsymbol{y} $。
\item $ d $ 是 \emph{对称的（symmetric）} ，即对所有 $ \boldsymbol{x}, \boldsymbol{y} \in V $，有 $ d(\boldsymbol{x}, \boldsymbol{y}) = d(\boldsymbol{y}, \boldsymbol{x}) $。
\item \emph{三角不等式（triangle inequality）} ：对所有 $ \boldsymbol{x}, \boldsymbol{y}, \boldsymbol{z} \in V $，有 $ d(\boldsymbol{x}, \boldsymbol{z}) \leqslant d(\boldsymbol{x}, \boldsymbol{y}) + d(\boldsymbol{y}, \boldsymbol{z}) $。
\end{enumerate}

\emph{备注} 。乍一看，内积和度量的性质列表非常相似。然而，通过对比定义 3.3 与定义 3.6，我们发现 $ \langle \boldsymbol{x}, \boldsymbol{y} \rangle $ 与 $ d(\boldsymbol{x}, \boldsymbol{y}) $ 的变化方向相反。非常相似的 $ \boldsymbol{x} $ 和 $ \boldsymbol{y} $ 会使内积取较大的值，而使度量取较小的值。♦
